% year: תשע"ז
% type: מבחן
% semester: א
% moed: א
% number: 2ב
% points: 10
% categories: func-limits, lhopital
% source: exams/מבחנים/תשעז/מועד א תשעז.pdf

\begin{question}
(10 נק') מצאו את הגבול הבא:
$\displaystyle \lim_{x\to 0}\frac{\sqrt[3]{8+3x}-2}{\sqrt[4]{16+5x}-2}$
\end{question}

\begin{solution}
בנקודה $x=0$: $\sqrt[3]{8}-2=0$ ו-$\sqrt[4]{16}-2=0$, כלומר גבול מהצורה $\frac00$. שתי הפונקציות גזירות בסביבת $0$:
\[
\left(\sqrt[3]{8+3x}\right)'=\frac13(8+3x)^{-2/3}\cdot3=(8+3x)^{-2/3},\qquad
\left(\sqrt[4]{16+5x}\right)'=\frac54(16+5x)^{-3/4},
\]
ונגזרת המכנה שונה מ-$0$ בסביבת $0$. לפי כלל לופיטל:
\[
\lim_{x\to0}\frac{\sqrt[3]{8+3x}-2}{\sqrt[4]{16+5x}-2}
=\lim_{x\to0}\frac{(8+3x)^{-2/3}}{\frac54(16+5x)^{-3/4}}
=\frac{8^{-2/3}}{\frac54\cdot16^{-3/4}}=\frac{\frac14}{\frac54\cdot\frac18}=\frac{\frac14}{\frac{5}{32}}=\frac85.
\]
\textbf{תשובה:} $\frac85$.
\end{solution}
